\section{Solutions for Chapter~\ref{ch:cosets}}

\paragraph{Exercise~\ref{ex:c-coset-1}: the perpendicular lattice.}
\phantomsection
\label{sol:ex:c-coset-1}
The orthogonal summands make the glue norm the sum
$2/3+2/3+1/2+78/36=4$. Its components have orders
$3,3,2,6$ modulo their respective lattices, so the full glue
class has order $\operatorname{lcm}(3,3,2,6)=6$.
Its pairings with the original lattice are integral: the
$A_2$ weights are dual-lattice vectors, while
$(\alpha/2,\alpha)=1$ and $(\gamma_{78}/6,\gamma_{78})=13$.
The even norm therefore gives an even integral index-six extension.
Its determinant is
\[
 \frac{3\cdot3\cdot2\cdot78}{6^2}
 =\frac{1404}{36}=39.
\]
If the plane lattice $L_0$ had a larger saturation of index
$r>1$ in $E_8$, that integral lattice would have determinant
$39/r^2$. Integrality of its Gram matrix makes this an integer,
so $r^2$ would divide the squarefree integer $39$, impossible.
This argument uses the actual integral ambient lattice; an
arbitrary real enlargement has no such determinant restriction.



\paragraph{Exercise~\ref{ex:c-coset-2}: why branching multiplicities square.}
\phantomsection
\label{sol:ex:c-coset-2}
For each lattice charge $n$, the three nonzero branching functions
are $\chi_n,\chi_n,2\chi_n$. Pairing each full multiplicity
space with its complex conjugate produces multiplicities
$1^2,1^2,2^2$. Summing gives $6Z_{L_0}$. The lattice vacuum
occurs at $n=0$, so the zero-energy multiplicity is six.
These are the multiplicities of the original embedded extension
$A=\one\oplus X\oplus2Y$. The calculation does not choose between
its possible $\mathcal H_2$ and $\mathcal H_3$ boundary categories,
nor replace it by a separately selected canonical algebra.

A partition function alone supplies no product on that
six-dimensional vacuum space. A full interpretation requires
a unital associative commutative vacuum product, its positive pairing and
adjoint, and compatible local OPE and sewing data for all other
sectors. In a unitary finite vacuum algebra the mutually orthogonal
primitive idempotents describe the separate vacuum components;
defect junctions must act consistently on this algebra. Six
disconnected identical lattice theories give one possible partition
function with this multiplicity, but do not automatically supply
the intended Haagerup action. A unique-vacuum, clustering theory
cannot simply keep the six-dimensional zero-energy space unchanged.



\paragraph{Exercise~\ref{ex:c-coset-3}: exact near-two charges.}
\phantomsection
\label{sol:ex:c-coset-3}
The affine charge is $3(40)/(40+2)=20/7$. Hence the diagonal
GKO charge is
\[
 1+\frac{117}{41}-\frac{20}{7}
 =\frac{286}{287}=1-\frac1{287}.
\]
Adding one Cartan stress gives $573/287=2-1/287$.
Its distance below two is $1/287$, approximately $.00348432$,
whereas $c_*-2$ is approximately $.025004294$ and is positive.
For completeness, the GKO stress by itself is $288/287$ below
two, and the diagonal affine stress is $6/7$ above two.
These exact rational stresses and the irrational axis are different
vectors and different chiral algebras; closeness of their charges
does not identify their full theories.



\paragraph{Exercise~\ref{ex:c-coset-4}: solving for a Virasoro axis.}
\phantomsection
\label{sol:ex:c-coset-4}
Expand $e*e$ using the table. The coefficient of $e_i$ is
$2a_i^2+t^2\lambda_i/f_i$ and that of $x$ is
$2t\lambda\cdot a$. Equating them to $2a_i$ and $2t$ gives
$a_i(1-a_i)=\lambda_i t^2/(2f_i)$ and the second equation.
The three ratios are
\[
 \frac{11/41}{286/287}=\frac7{26},\qquad
 \frac{192/1271}{2400/1271}=\frac2{25},\qquad
 \frac{49/31}{70/31}=\frac7{10},
\]
giving the stated $q_i$.
Let $\omega'=e_0+e_1+e_2$. Invariance and $\omega'*e=2e$
give
$2\langle e,e\rangle=\langle\omega',e*e\rangle
=2\langle\omega',e\rangle$. Therefore
$c(e)=2\langle e,e\rangle=2\sum_i f_i a_i$ on the
idempotent locus.



\paragraph{Exercise~\ref{ex:c-coset-5}: the vacuum grade-four Gram matrix.}
\phantomsection
\label{sol:ex:c-coset-5}
Using $[L_m,L_n]=(m-n)L_{m+n}+\frac c{12}(m^3-m)\delta_{m+n,0}$
and $L_{-1}\one=0$, the matrix in the stated basis is
\[
 \begin{pmatrix}5c&3c\\3c&c(c+8)/2\end{pmatrix}.
\]
For example the mixed entry follows from
$L_4L_{-2}^2\one=6L_2L_{-2}\one=3c\one$.
The first diagonal entry is
$\langle\one,[L_4,L_{-4}]\one\rangle=5c$.
For the other one, commute $L_2$ past both factors to obtain
$L_2L_{-2}^2\one=(8+c)L_{-2}\one$; a second $L_2$ gives
$(8+c)c/2$. These are inner products because the unitary
adjoint is $L_n^\dagger=L_{-n}$ and $\|\one\|=1$.
Its determinant is
$5c^2(c+8)/2-9c^2=c^2(5c+22)/2$, positive for $c>0$.
The two vacuum descendants are therefore independent at $c_*$.



\paragraph{Exercise~\ref{ex:c-coset-6}: the trivial grading.}
\phantomsection
\label{sol:ex:c-coset-6}
Let $g_X$ be the group element assigned to a simple $X$, with
$g_Z=g_Xg_Y$ whenever $Z$ occurs in $X\otimes Y$.
The tensor unit has $g_{\one}=1$, either by the unital assignment
or by $g_{\one}^2=g_{\one}$ in a group.
Since both $\one$ and $\rho$ occur in $\rho^2$,
$g_\rho^2=1=g_\rho$. Thus $g_\rho=1$.
The constituent $\alpha\rho$ then implies
$g_\alpha g_\rho=g_{\alpha\rho}=g_\rho^2=1$,
so $g_\alpha=1$; the same holds for every remaining simple.
The group need not be abelian for this argument.

To apply this to currents, each simple line must transmit the
chosen current algebra by a definite invertible rotation, and
every nonzero fusion junction must identify the rotation of
each output constituent with the product of the two input
rotations. Those rotations then give the required group-valued
assignment. Preserving only the stress, or retaining arbitrary
matrix elements of a line after a projection, does not establish
this transmission property.



\paragraph{Exercise~\ref{ex:c-coset-7}: bounded projected energy.}
\phantomsection
\label{sol:ex:c-coset-7}
Project $(m,n)\in\Z^2$ to $m+\theta n$ with irrational $\theta$.
There are infinitely many pairs whose projected value lies in any
fixed sufficiently large interval: for each $n$, choose $m$ nearest
to $-\theta n$, giving absolute value at most $1/2$. Their full
Euclidean lengths grow, so a finite ambient-energy cutoff would
remove all but finitely many. Discarding the perpendicular energy
removes that control. To put this in the form of positive commuting
energies, use
\[
 E_{\parallel}=\frac{(m+\theta n)^2}{1+\theta^2},\qquad
 E_{\perp}=\frac{(n-\theta m)^2}{1+\theta^2},\qquad
 E_{\parallel}+E_{\perp}=m^2+n^2.
\]
The selected states have $E_{\parallel}\leq1/[4(1+\theta^2)]$,
so their reduced thermal sum already diverges at any positive
temperature. Irrationality makes the linear projection injective:
the issue is infinitely many states in a bounded energy interval,
not necessarily infinite multiplicity at one exact eigenvalue.
This example illustrates the convergence
problem; it does not identify the actual irrational commutant
branching with a freely projected lattice.



\paragraph{Exercise~\ref{ex:c-coset-8}: a sum and its diagonal square.}
\phantomsection
\label{sol:ex:c-coset-8}
The example $a_n=1/n$ has a divergent sum and a convergent square
sum. Therefore the divergence of $R(\beta)$ alone would be
insufficient. A character is more special: write it as a sum over
an actual basis,
$\chi_\lambda(\beta)=\sum_j e^{-\beta E_j}$, including repetitions
for multiplicities. Its square is a sum over ordered pairs of basis
vectors. Keeping only the pairs $(j,j)$ gives
$\chi_\lambda(\beta)^2\geq\sum_j e^{-2\beta E_j}
=\chi_\lambda(2\beta)$.
All terms are positive, and integer multiplicities permit this
basis interpretation: a multiplicity $d$ contributes $d^2$
ordered pairs and at least $d$ diagonal pairs. The shifted energies
$E_j$ may be negative; positivity of their exponentials is all
that is used. Since $R(2\beta)$ diverges for every
$\beta>0$, the diagonal thermal sum diverges too.



\paragraph{Exercise~\ref{ex:c-coset-9}: independent temperatures.}
\phantomsection
\label{sol:ex:c-coset-9}
Write $s=\beta+\bar\beta$, $p=\beta\bar\beta$, and
$t=4\pi^2s/p$. The inverted temperatures sum to $t$, while
eta inversion gives
\[
 \eta(4\pi^2/\beta)\eta(4\pi^2/\bar\beta)
 =\frac{\sqrt p}{2\pi}\eta(\beta)\eta(\bar\beta).
\]
Equating the original and inverted amplitudes therefore gives
$G(t)=\sqrt p\,G(s)/(2\pi)=\sqrt{s/t}\,G(s)$.
For fixed $s>0$, independently positive temperatures allow
all $0<p\leq s^2/4$, equivalently all $t$ with
$st\geq16\pi^2$. Given $s_1,s_2>0$, choose
$u\geq16\pi^2\max(s_1^{-1},s_2^{-1})$. Both pairs $(s_i,u)$
are allowed, forcing $\sqrt{s_1}G(s_1)=\sqrt{s_2}G(s_2)$.

If only equal temperatures were checked, one would have
$p=s^2/4$ and only the involution $t=16\pi^2/s$.
It relates each argument to one partner, with no freedom to
connect two arbitrary arguments through a common $u$.
An invariance under this involution alone does not force
$\sqrt{s}G(s)$ to be constant. The stronger two-variable
modularity hypothesis is essential to the uniqueness proof.


\paragraph{Exercise~\ref{ex:c-coset-10}: the dimension threshold.}
\phantomsection
\label{sol:ex:c-coset-10}
For a noninvertible simple $X$, the product $X\otimes X^*$ contains
the unit and at least one further simple. All simple
Frobenius--Perron dimensions are at least one, so
$\FPdim(X)^2\geq2$. If there are at least three simples, two
others each contribute at least one to global dimension, proving
the bound four.

At rank two, the nonunit simple is self-dual and
$X^2=\one+nX$ for an integer $n\geq1$. Its positive dimension
is $(n+\sqrt{n^2+4})/2$, an increasing function of $n$. At $n=1$
the global dimension is $1+((1+\sqrt5)/2)^2=(5+\sqrt5)/2$.
In the descent argument, a nontrivial central adjoint algebra would
have dimension at least $d^2$, forcing the image category below
both nonpointed thresholds:
\[
 \operatorname{Dim}\mathcal D
 \leq\frac{3(1+d^2)}{d^2}=3+\frac3{d^2}
 <\frac{10}{3}<\frac{5+\sqrt5}{2},
\]
where $d>3$ gives the strict bound. Its being pointed would
contradict the irrational dimension $d$ of the image of $\rho$,
since any object of a pointed category is a direct sum of
dimension-one simples. The adjoint algebra in this argument is
$R(\one)$ for the hypothesized descent functor; it is not the
original $E_8/U$ extension algebra, despite the chapter's reused
local letter $A$.

